% FIND -- os/samir/cmd/query.c q_find (unquoted literal through ndx_seek on
% the master index; SET EXACT; #26 without an index).  Gate: test-ndx-seek,
% test-samir-query.
\begin{XbaseEntry}{FIND}{command}
Finds the first record whose index key begins with a literal character
string.
\Syntax
\SyntaxLine{\kw{FIND} \xp{literal string}}
\Usage
\cmd{FIND} searches the master index of the database in USE (the first index
named in \cmd{USE}, or the one chosen with SET ORDER) for the characters
typed after the command, without quotation marks.  If a key begins with them,
the record pointer moves to the first such record in index order and
\fn{FOUND()} is true.  If none does, the pointer moves to the end of the file:
\fn{EOF()} is true and \fn{FOUND()} false.

While SET EXACT is OFF, a key that begins with the string is found; with SET
EXACT ON, the key must equal the string, apart from trailing blanks.  Capital
and small letters are different.

\cmd{FIND} needs an open index.  Without one it displays
\xmsg{26}{Database is not indexed.}
\Tips
\cmd{SEEK} does the same with an expression, so it can search for the
contents of a memory variable or for a number.
\Examples
\DotExample{examples/c-find}
A string that begins no key leaves the pointer at the end of the file:
\DotExample{examples/c-find-2}
\SeeAlso{SEEK, LOCATE, SET EXACT, SET ORDER, USE, FOUND()}
\end{XbaseEntry}

